\hypertarget{appendix-ed3f9bc298fe}{}\label{referencestart:appendix-ed3f9bc298fe}
\section{Triangle rasterization and coverage tests}\label{appendix-ed3f9bc298fe-triangle-rasterization-and-coverage-tests}

\textbf{Curriculum depth:} Developed · \textbf{First seen:} University courses · \textbf{Latest development:} University courses

\subsection{Definition and mechanism}\label{appendix-ed3f9bc298fe-definition-and-mechanism}

Rasterization converts triangle geometry into a finite set of image samples. A pixel's ideal coverage is the area fraction inside the triangle; a center test approximates this integral with one binary answer and therefore creates jagged edges. A directed edge defines an inside half-plane, so three signed edge tests classify a point against a triangle. A consistent tie rule assigns samples exactly on shared edges to one triangle, preventing cracks and double coverage. Supersampling tests several locations per pixel and averages their coverage, trading cost for a better area estimate. Hardware often checks candidate samples in parallel within a bounding box, then rejects or accepts coarse blocks that are wholly outside or inside. Triangle coverage is also the basis for interpolating depth, color, normals, and texture coordinates.

\subsection{Prerequisites}\label{appendix-ed3f9bc298fe-prerequisites}

\begin{itemize}
\tightlist
\item
  \hyperref[appendix-02b9d9a81dca]{Barycentric coordinates and interpolation}
\item
  \hyperref[appendix-6a7869dd79a6]{Sampling aliasing and reconstruction}
\end{itemize}

\subsection{Used by}\label{appendix-ed3f9bc298fe-used-by}

\begin{itemize}
\item
  {Worlds, Images and Inference}
\item
  \hyperref[appendix-97578ed97952]{Depth buffering and visibility order}
\item
  \hyperref[appendix-40426e4b2103]{Homogeneous clipping and perspective interpolation} \#\# Full teaching explanation
\end{itemize}

\href{../03-Geometry-Cameras-and-Image-Formation/main.pdf\#nameddest=03-visual-computing-i-ch001-a-scene-description-becomes-an-image}{A scene description becomes an image (Geometry, Cameras, and Image Formation, Geometry, Cameras, and Image Formation, Chapter 8)} · \href{../02-Mathematics-for-Visual-Computing/main.pdf\#nameddest=03-visual-computing-i-ch012-rasterization-coverage-and-depth}{Rasterization, coverage, and depth (Mathematics for Visual Computing, Mathematics for Visual Computing, Chapter 16)}
